\section{Causal structure: the six-dimensional baseline}
\label{sec:baseline}

Every construction so far has used a Lorentzian four-plane chosen inside $\R^{3,3}$: the interior plane $E$ of the Dirac structure, or the slice $\Sigma$
of the curved model. This part asks whether that choice, and with it the light cone and the distinction between past and future, can be obtained from the
split geometry rather than postulated. Two levels must be distinguished. \emph{Weak emergence} means that an already selected $3{+}1$ sector carries a
consistent quantum causal structure. \emph{Strong emergence} means that the dynamics determine why that sector, with that causal structure, arises from the
six-dimensional theory. The light cone must not be put into the admissibility conditions and then be said to emerge. We fix the failure conditions in advance:
uncontrolled growing modes in the six-dimensional theory; a projector onto the four-dimensional sector imposed by hand; ghosts or negative spectral weights in
the effective theory; propagation outside the derived cone; extra modes or imposed equations in the Maxwell sector.

\subsection{The unrestricted scalar}

\begin{proposition}[A free field on $\R^{3,3}$]
\label{prop:toy}
For the action $S_6=\tfrac12\int d^6X\,(\eta^{AB}\partial_A\phi\,\partial_B\phi-m^2\phi^2)$ in coordinates $(t,u,v,x,y,z)$, with $t,u,v$ temporal:
\begin{enumerate}[label=(\roman*),leftmargin=2.2em]
\item plane waves $e^{-i\omega t+i(q_1u+q_2v)+i\bk\cdot\mathbf x}$ obey
\begin{equation}
\label{eq:toydisp}
\omega^2=|\bk|^2+m^2-q^2,\qquad q^2=q_1^2+q_2^2 ;
\end{equation}
\item for $q^2>|\bk|^2+m^2$ the growth rate is $\gamma=\sqrt{q^2-|\bk|^2-m^2}$, and $(1+q^2)^{-s/2}e^\gamma\to\infty$ for every $s$: data small in any Sobolev norm
produce arbitrarily large solutions at $t=1$, so the Cauchy problem with respect to $t$ is ill-posed in the sense of Hadamard \cite{Hadamard23};
\item the canonical energy density $\tfrac12(\pi^2-\phi_u^2-\phi_v^2+|\nabla_{\mathbf x}\phi|^2+m^2\phi^2)$ averages to $\tfrac14a^2(m^2-q^2)$ on $\phi=a\sin qu$: the
energy is unbounded below.
\end{enumerate}
\end{proposition}

\begin{figure}[t]
\centering
\begin{tikzpicture}
\begin{axis}[width=10.5cm,height=6.6cm,xmin=0,xmax=3.2,ymin=-5.5,ymax=4.5,axis on top,clip=false,
  axis lines=left,
  xlabel={surplus momentum $q$},ylabel={$\omega^2$},xtick={1.4142},xticklabels={$\sqrt{|\bk|^2+m^2}$},ytick={0},
  ylabel style={rotate=-90}]
\addplot[draw=none,fill=scol!12,forget plot] coordinates {(1.4142,-5.5) (3.2,-5.5) (3.2,4.5) (1.4142,4.5)} -- cycle;
\draw[gray,dashed] (axis cs:0,0) -- (axis cs:3.2,0);
\addplot[tcol,very thick,domain=0:2.739,samples=80] {2-x^2};
\addplot[only marks,mark=*,mark size=2.6pt,gcol] coordinates {(0,2)};
\node[gcol,anchor=south west,font=\small,align=left] at (axis cs:0.06,2.2) {$q=0$: the four-dimensional\\ mass shell, kept by the constraints};
\node[scol,anchor=north west,font=\small,align=left] at (axis cs:2.25,4.2) {unstable band:\\ $\omega=\pm i\sqrt{q^2-|\bk|^2-m^2}$};
\end{axis}
\end{tikzpicture}
\caption{The dispersion relation of a free scalar on $\R^{3,3}$ at fixed spatial momentum (\cref{prop:toy}). Every surplus momentum beyond $\sqrt{|\bk|^2+m^2}$ gives
an exponentially growing mode, with a rate that grows without bound. The constraints of \cref{prop:constraints} keep exactly the point $q=0$.}
\label{fig:toy}
\end{figure}

\begin{proposition}[No time orientation in signature $(3,3)$]
\label{prop:orient}
The timelike set $\{\eta(X,X)>0\}$ of $\R^{p,q}$ is connected for $p\ge2$: the path $X(s)=(\cos s,\sin s,0;\,0,0,0)$ joins $X$ to $-X$ with $\eta=1$ throughout. Only for $p=1$
does it split into a past and a future component.
\end{proposition}

Causal order is therefore not merely unproven in $\R^{3,3}$; it is absent. Anything called a light cone with a future has to come from a reduction to a single
time. The arena already contains such a reduction, as an input. For the geodesic test particle of \cref{prop:constraints}, $H=\tfrac12(h^{\mu\nu}k_\mu k_\nu+p_u^2+p_v^2)$
with $h=\operatorname{diag}(1,-1,-1,-1)$, the mass shell $2H=m^2$ is $p_0^2=|\mathbf p|^2+m^2-p_u^2-p_v^2$, the relation \eqref{eq:toydisp} with $q=(p_u,p_v)$. The
first-class constraints $p_u=p_v=0$ remove exactly the unstable band, as \cref{fig:toy} shows.

\begin{example}[Compact surplus times give tachyons]
\label{ex:tachyons}
If instead $u$ and $v$ are made periodic with period $2\pi R$, the surplus momenta become $n/R$ with $n\in\Z^2$, and the four-dimensional masses are
\[
m_n^2=m^2-\frac{|n|^2}{R^2}.
\]
Every mode with $|n|>mR$ is a tachyon, and there are infinitely many: for $mR=3$ there are $92$, $412$ and $1652$ of them within $|n_i|\le5$, $10$ and $20$.
Compactification, the usual cure for extra space dimensions, is no cure for extra times.
\end{example}

An ultrahyperbolic problem can be made well posed by a nonlocal constraint on the data \cite{CraigWeinstein09}, and the dimensionality of spacetime has been
argued from such considerations \cite{Tegmark97}. The constraints of the arena, $\partial_u\phi=\partial_v\phi=0$, are local and stronger: they keep only the zero modes,
which is dimensional reduction. Either way the cure is a restriction, and at this stage it is an input. \Cref{sec:constraints} asks where it can come from; the next
two sections ask whether the four-plane it presupposes can be selected.
